%% Dedicatória, agradecimentos e epígrafe.
%% O texto original não tem nenhum dos três; descomente e preencha
%% os que quiser incluir.

% Dedicatória (poucas linhas):
%\dedicatoria{Texto da dedicatória.}

% Agradecimentos (uma única página), texto do autor. Os parágrafos
% ficam separados por um pequeno espaço e, a partir do segundo, com o
% recuo do corpo do texto.
\prefacesection{Agradecimentos}
\begingroup\setlength{\parskip}{\medskipamount}
A Deus, por me sustentar ao longo destes anos.

Ao meu orientador, Prof. Dr. Eduardo Paiva Okabe, pela confiança depositada em minha formação acadêmica desde a segunda disciplina em que assisti às suas aulas, Cinemática de Mecanismos, base da teoria aplicada neste trabalho, e por me aceitar como seu monitor na mesma disciplina no ano seguinte. Agradeço por ter me proposto um projeto de Iniciação Científica na mesma área e pela oportunidade de colaborar no artigo \textit{Modeling and simulation of an upper limb exoskeleton}, publicado em periódico internacional. Neste trabalho, agradeço a paciência, a orientação e a validação ao longo dos últimos meses, e o planejamento dos próximos passos em uma janela de tempo muito curta. Às quartas-feiras, depois das 17h30, e nas reuniões online.

Ao Prof. Dr. José Luiz Pereira Brittes, que, além das lições aprendidas em sala de aula, me ensinou a noção e a necessidade da educação continuada como um ato de amor às exatas, acompanhada de uma severa crítica ao utilitarismo que tantas vezes se prega nos circuitos em que vivemos. A ele agradeço também por identificar minha afinidade com os temas acadêmicos antes de mim, por acreditar no projeto de Exoesqueletos para assistência de pessoas com perdas neuromusculares e me envolver nele, e pela assistência em minha parte no artigo já citado.

Ao Prof. Dr. Felipe Wenzel da Silva Tuckmantel, por ter aceitado o convite para compor a banca examinadora e por somar seus conhecimentos de pesquisa ao tema deste trabalho.

Aos meus pais, Valter Luís de Brito e Carla Simone de Almeida Brito, professores de formação. Por acreditarem na universidade pública. E por terem sido, em meio a desabafos e descrenças, um notável alento.

Às agências de fomento ao ensino e à pesquisa, pela oportunidade de ser monitor em quatro ocasiões e de realizar dois projetos de Iniciação Científica.

Ao meu irmão, Daniel Vinícius de Almeida Brito, força motriz do esforço que me permitiu abrir caminho para que ele também estudasse, sem onerar o orçamento familiar.

Aos muitos colegas e amigos, pela convivência e pelo companheirismo ao longo do curso, em especial aos grupos Confraria FCA e Alvinegros FCA, formados entre os ingressantes de 2022 dos cursos de Engenharia de Manufatura e Engenharia de Produção.

Às comunidades que desenvolvem e mantêm o Python e as bibliotecas científicas livres, empregados neste trabalho.
\par\endgroup

% Epígrafe:
%\epigrafe{%
%Texto da epígrafe.
%
%  \begin{flushright}
%    Autor da citação
%  \end{flushright}
%}
